\chapter{The project in plain language}\label{ch:summary}
This chapter explains, without equations, what was done, what the model is and how well it explains the papers. Every later chapter expands one part of it.

\section{What we were trying to do}
Manas Geeta Arun and colleagues measured how a fruit fly's genes affect fitness in males versus females. The key number is the \textbf{intersexual genetic correlation for fitness}, $r$:
\begin{itemize}
\item $r>0$: genomes that are good for males are also good for females. Most variation is ``shared'', e.g.\ harmful mutations that hurt both sexes.
\item $r<0$: genomes good for one sex are bad for the other. This is \textbf{sexual antagonism (SA)}: the same genes pull the sexes in opposite directions.
\end{itemize}
Manas found $r\approx+0.25$ to $+0.40$ in his population (LH). Earlier studies on the parent population (LH$_\mathrm{M}$) found $r\approx-0.30$ to $-0.52$. His other papers argue two things:
\begin{enumerate}
\item a negative $r$ can appear even without true SA, through \textbf{genetic linkage} (Hill--Robertson interference, HRI);
\item ``modifier'' genes can resolve the conflict, more easily on the X chromosome.
\end{enumerate}
The job was to build a simulator that explains those results and reconciles the conflicting ones, \emph{without} simply being tuned to reproduce them.

\section{What the model is: three pieces}
\begin{enumerate}
\item \textbf{The maths layer (exact equations).} The classical population-genetic formulas: how allele frequencies change when selection differs between the sexes, with linkage between genes. For the modifier paper we re-ran the authors' own code on their exact parameter grid.
\item \textbf{A virtual fly population (the core).} The computer simulates thousands of individual flies over thousands of generations.
  \begin{itemize}
  \item Each fly has a genome of about 40,000 possible mutation sites on one autosome plus the X.
  \item New mutations appear each generation. Each is one of four types:
    \begin{itemize}
    \item female-only harmful;
    \item male-only harmful;
    \item shared harmful (hurts both sexes);
    \item antagonistic (good for one sex, bad for the other).
    \end{itemize}
  \item Real \emph{Drosophila} rules apply: males have one X; males do not recombine their chromosomes, females do.
  \item Fitter flies leave more offspring, so selection, drift and linkage happen \emph{naturally}.
  \end{itemize}
  The simulator never sees the equations; the patterns have to emerge.
\item \textbf{A virtual copy of Manas's experiment.} Inside the simulated population we repeat his lab protocol:
  \begin{itemize}
  \item sample 39 genomes (``hemigenomes'') the way the clone-generator flies do;
  \item put each into males and into females;
  \item measure eggs laid and offspring sired in vials at the three sex ratios;
  \item run his exact statistics on the simulated data.
  \end{itemize}
\end{enumerate}
\begin{keybox}[The honesty rule]
Real measurements are noisy, so the virtual assay needs a noise level. It was tuned \emph{only} to how variable his raw data were, never to the correlations. Everything else is a genuine prediction: $r$, the sex-ratio pattern and the other correlations.
\end{keybox}

\section{How well it explains each result}
\begin{longtable}{p{3.3cm}p{4.3cm}p{7.6cm}}\toprule
Paper & What they found & How the model does\\\midrule\endhead
Sex-ratio study (BMC 2022) & $r=+0.38/+0.40/+0.25$ at male-biased / equal / female-biased sex ratios; the difference is not significant & \textbf{Recomputed exactly} from his raw data, to 4 decimals. This works only \emph{without} a data transformation his Methods describe, a mismatch we found along the way.

\textbf{Simulated:} a population with mostly shared mutations and few antagonistic ones (about 6\% of new mutations) gives $r\approx+0.45$. \emph{Audit:} that architecture was \emph{selected} from a sweep because it matched, so this is a fit, not a prediction (Chapter~\ref{ch:audit}).\\
 & The sex-ratio difference & Explained as \emph{measurement dilution}: at the female-biased ratio the genetic signal is weaker relative to noise, so $r$ looks smaller although the genes have not changed. \emph{Audit:} the same ordering follows from the raw data's reliability alone, with no simulator. Attenuation explains about a third of the drop; equal ``true'' $r$ is not rejected ($p\approx0.3$), but the test has little power.\\
 & Other correlations & Male across-sex-ratio correlations agree (0.51 / 0.68 / 0.49 vs 0.56 / 0.70 / 0.54), but the noise level was calibrated on the same data. The female ones are too high in the model (tension, $p=0.017$), which suggests a real genotype $\times$ sex-ratio effect.\\
HRI preprint (2025) & Linkage alone can make $r$ negative (about $-0.08$), more so with tight linkage and more mutation & \textbf{Reproduced:} $r$ is negative, weakens as recombination increases, is about 0 for loosely linked genes, and strengthens with mutation rate. At their exact settings: $-0.04$ vs their $-0.08$.\\
Modifier preprint (2023) & Perfect modifiers resolve the conflict in about 96\% of cases; the X is better than autosomes & \textbf{6 of 7 claims reproduced} by re-running their code (97\% vs 96\%). One claim (autosomes beat the X at very high modifier strength) did not hold.\\
Am Nat 2025 & -- & \textbf{Could not assess}; the paper was inaccessible.\\
Conflicting LH$_\mathrm{M}$ studies & $r=-0.3$ to $-0.5$ & \textbf{Reconciled, not fitted.} Raising the antagonistic share from about 6\% to 10--19\% of new mutations moves the simulated $r$ from $+0.45$ to $-0.25$ to $-0.35$. Neither sampling noise nor linkage alone can make that jump: LH and LH$_\mathrm{M}$ most likely differ in their mutation mix or population history.\\\bottomrule
\end{longtable}

\section{What we learned that is new}
\begin{itemize}
\item \textbf{Small changes, big effect.} A few percent more antagonistic mutations flips the sign of $r$. Populations with identical genetic make-up can end at very different $r$ by chance history.
\item \textbf{Linkage alone cannot explain big negative $r$.} HRI gives only small negative values (about $-0.18$ at most on average, at the tightest linkage) and can never make $r$ positive. Manas's positive $r$ is therefore not an HRI artefact, and the strongly negative LH$_\mathrm{M}$ values probably need real antagonism.
\item \textbf{A methods warning.} Male flies do not recombine, so the sexes can become genetically different. The HRI paper's way of computing $r$, pooling male and female genotypes, then gives a misleadingly negative number. The hemiclone method does not have this problem.
\item \textbf{The ``SA proportion'' adds nothing beyond $r$.} For standardised line means it equals $(1-r)/2$ exactly (\S\ref{ex:rot}). The published proportions are this identity applied to the published $r$.
\item \textbf{The obvious test fails.} Letting the genomes recombine in the lab, to break linkage and see whether $r$ disappears, does not work in simulation: the linkage that matters is too tight to break in a few generations. Comparing small and large populations works in principle, but needs about 30 replicate populations per size.
\end{itemize}

\section{The honest limits}
\begin{itemize}
\item The model shows which genetic make-ups \emph{could} produce each result. It cannot measure which one LH actually has: nobody has measured the mutation mix.
\item Sex ratio enters only as ``stronger or weaker selection''. With no explicit harm or resistance traits, the model cannot say whether the sex-ratio effect is dilution or a real change in which genes matter. 39 lines cannot tell these apart either.
\item To keep run times manageable, smaller populations were scaled to mimic large ones. This preserves direction and ordering but exaggerates linkage effects about 2$\times$ in size. The most important case was checked at full scale.
\item The machine-learning summary of the simulations (a Gaussian process) failed its accuracy check, so nothing relies on it.
\end{itemize}

\section{Bottom line (after the audit, Chapter~\ref{ch:audit})}
\begin{itemize}
\item \textbf{Genuine predictions agree in direction}, with weak to moderate evidence:
  \begin{itemize}
  \item linkage-driven negative $r$ that fades with recombination (a replication of the preprint's own model);
  \item Khan's preference for healthy females, predicted from an independent calibration (effect about half to two-thirds of the observed size);
  \item the courts-most spread in held-out blocks.
  \end{itemize}
\item \textbf{Several apparent matches are not evidence.}
  \begin{itemize}
  \item The LH value of $+0.45$ was selected.
  \item The sex-ratio ordering is an echo of the variance calibration.
  \item The ``SA proportion'' is an identity.
  \item Chinmay's courtship and telegony nulls follow from the assumptions.
  \end{itemize}
\item \textbf{The model is rejected or in tension} where the data contain structure it lacks: Khan's between-vial dispersion, the telegony sex effect, genotype effects on latency, and female genotype $\times$ sex-ratio interaction.
\item The conflict between labs is \emph{compatible} with a modest difference in antagonistic mutational input or population history. Sampling noise and linkage alone cannot produce it. This is an exploratory statement, not a fitted explanation.
\item The model \textbf{cannot} pin down LH's true genetic make-up, decide whether sex ratio changes which genes matter, or say why Chinmay's males showed no choice. Those need new data: above all, fecundity of the choice-vial females, and a replicated telegony block.
\end{itemize}
